% Tiny contract
% Teams: Navigation
% Requirements: REQ-OPS-050, REQ-FUN-060, REQ-FUN-240
% Status: done

The navigation subsystem performs autonomous path planning, trajectory tracking, and collision avoidance. After receiving a goal pose or waypoint sequence, the rover executes motion in a closed loop: sensor updates (\autoref{par:navigation_sensor_suite}), costmap update, global planning, local trajectory tracking, and command execution. The same stack is used for traversal and task positioning, while manipulation is handled separately (\autoref{par:arm_kinematics_motion_planning}).

Global planning uses NavFn (Dijkstra) over a \qty{0.05}{\metre} resolution 2D costmap. Local motion tracks the global path while reacting in real time to terrain changes and newly detected obstacles. Collision avoidance, including other rovers (REQ-FUN-240), is enforced through obstacle detection and safety-margin inflation.

The costmap integrates depth and range sensing into voxel, obstacle, and inflation layers at \qty{0.05}{\metre} resolution. Obstacles include occupied cells, terrain protrusions above \qty{0.15}{\metre}, depressions deeper than \qty{0.10}{\metre}, and cells with cost $\geq 253$. Traversability is estimated from terrain geometry, with safe regions defined by cost thresholds.

If the path becomes obstructed, the navigation state machine (Figure~\ref{fig:navigation_combined}) performs local recovery, triggers global replanning if needed, or executes a fail-safe stop if no safe trajectory is found within a timeout.
\begin{figure}[H]
  \centering

  \begin{minipage}{0.49\linewidth}
    \centering
    \resizebox{\linewidth}{!}{% Nav2 dataflow diagram - corrected final version
% Requires in preamble.tex:
% \usetikzlibrary{arrows.meta,positioning,calc,backgrounds,fit}

\tikzset{
  blk/.style={
    rectangle, draw=black, fill=white, text=black, line width=0.5pt,
    minimum width=4.2cm, minimum height=0.95cm,
    align=center, font=\small,
    inner ysep=2pt
  },
  arr/.style={-{Stealth[length=2.5pt]}, line width=0.5pt, draw=black},
}

\begin{tikzpicture}[node distance=1.45cm and 2.55cm]

  %% === SENSORS COLUMN ===
  \node[blk] (stereo) {StereoLabs ZED 2i\\[-1pt]{\scriptsize USB 3.0 - 110\textdegree$\times$70\textdegree{} FOV}};
  \node[blk] (tof) [below=of stereo] {ToF Cameras $\times$3\\[-1pt]{\scriptsize MaixSense A010 - I2C}};
  \node[blk] (lidar) [below=of tof] {RPLIDAR-S2\\[-1pt]{\scriptsize UART - 360\textdegree{} - 10\,Hz}};

  %% === COSTMAP COLUMN ===
  \node[blk] (voxel) [right=of stereo] {Voxel Layer\\[-1pt]{\scriptsize 3D grid - 0.05\,m resolution}};
  \node[blk] (obstacle) [below=of voxel] {Obstacle Layer\\[-1pt]{\scriptsize 2D projection}};
  \node[blk] (inflation) [below=of obstacle] {Inflation Layer\\[-1pt]{\scriptsize radius 0.15\,m}};

  %% === NAV2 COLUMN ===
  \node[blk] (global) [right=3.45cm of voxel] {NavFn Planner\\[-1pt]{\scriptsize Dijkstra - $<500$\,ms}};
  \node[blk] (local) [below=of global] {DWB Controller\\[-1pt]{\scriptsize 20\,Hz - $(v,\omega)$}};
  \node[blk] (cmdvel) [below=of local] {Motor Controllers};

  %% === SENSOR INPUTS ===
  \draw[arr] (stereo.east) -- ++(0.80,0) |- (voxel.west);
  \node[font=\scriptsize] at ($(stereo.east)!0.56!(voxel.west)+(0,0.28)$) {Depth};

  \draw[arr] (tof.east) -- ++(0.80,0) |- (voxel.west);
  \node[font=\scriptsize] at ($(tof.east)!0.56!(voxel.west)+(0,0.75)$) {Range};

  \draw[arr] (lidar.east) -- ++(1.90,0) |- (obstacle.west);
  \node[font=\scriptsize] at ($(lidar.east)+(1.10,0.28)$) {LaserScan};

  %% === COSTMAP PIPELINE ===
  \draw[arr] (voxel.south) -- (obstacle.north)
    node[midway, right, font=\scriptsize] {voxels};

  \draw[arr] (obstacle.south) -- (inflation.north)
    node[midway, right, font=\scriptsize] {obstacles};

  %% === COSTMAP -> GLOBAL PLANNER ===
  \coordinate (cgA) at ($(inflation.east)+(0.60,0)$);
  \coordinate (cgB) at ($(cgA |- global.west)$);
  \draw[arr] (inflation.east) -- (cgA) -- (cgB) -- (global.west);
  \node[font=\scriptsize, above] at ($(cgB)!0.5!(global.west)+(0.00,0.12)$) {Costmap};

  %% === COSTMAP -> LOCAL CONTROLLER ===
  \coordinate (clA) at ($(inflation.east)+(1.05,0)$);
  \coordinate (clB) at ($(clA |- local.west)$);
  \draw[arr] (inflation.east) -- (clA) -- (clB) -- (local.west);
  \node[font=\scriptsize, above] at ($(clB)!0.5!(local.west)+(-0.25,0.12)$) {Local costmap};

  %% === GLOBAL PLAN -> LOCAL CONTROLLER ===
  \draw[arr] (global.south) -- (local.north);
  \node[font=\scriptsize, right, align=left]
    at ($(global.south)!0.5!(local.north)+(0.16,0.00)$)
    {\texttt{nav\_msgs/}\\\texttt{Path}};

  %% === OUTPUT ===
  \draw[arr] (local.south) -- (cmdvel.north);
  \node[font=\scriptsize, right, align=left]
    at ($(local.south)!0.56!(cmdvel.north)+(0.16,0.00)$)
    {\texttt{geometry\_msgs/}\\\texttt{Twist}};

  %% === COLUMN LABELS ===
  \node[font=\small\bfseries, above=0.55cm of stereo] {SENSOR SUITE};
  \node[font=\small\bfseries, above=0.55cm of voxel] {COSTMAP};
  \node[font=\small\bfseries, above=0.55cm of global] {NAV2 STACK};

  %% === SECTION FRAMES ===
  \begin{scope}[on background layer]
    \node[
      draw=black, fill=black!5, dashed, line width=0.8pt,
      inner sep=8pt, rounded corners=3pt,
      fit=(stereo)(lidar)
    ] {};

    \node[
      draw=black, fill=black!5, dashed, line width=0.8pt,
      inner sep=8pt, rounded corners=3pt,
      fit=(voxel)(inflation)
    ] {};

    \node[
      draw=ercblue, fill=ercblue!5, dashed, line width=0.8pt,
      inner sep=8pt, rounded corners=3pt,
      fit=(global)(cmdvel)
    ] {};
  \end{scope}

\end{tikzpicture}}
  \end{minipage}
  \hfill
  \begin{minipage}{0.49\linewidth}
    \centering
    \resizebox{\linewidth}{!}{% Navigation state machine - corrected with larger label spacing
% Requires:
% \usetikzlibrary{arrows.meta,positioning,calc}

\tikzset{
  blk/.style={
    rectangle, draw=black, fill=white, line width=0.5pt,
    minimum width=3.8cm, minimum height=0.9cm,
    align=center, font=\small
  },
  arr/.style={-{Stealth[length=2.5pt]}, line width=0.5pt, draw=black},
  dlink/.style={-{Stealth[length=2.5pt]}, line width=0.5pt, dashed, draw=black},
}

\begin{tikzpicture}[node distance=1.9cm and 2.5cm]

  %% === STATES - Row 1 ===
  \node[blk] (idle) {IDLE\\{\scriptsize waiting for goal}};
  \node[blk] (compute) [right=3.0cm of idle]
    {COMPUTE PATH\\{\scriptsize NavFn - Dijkstra}};
  \node[blk] (drive) [right=of compute]
    {TRACK TRAJECTORY\\{\scriptsize DWB - 20\,Hz}};

  %% === STATES - Row 2 ===
  \node[blk] (waypoint) [below=of idle]
    {WAYPOINT REACHED\\{\scriptsize tolerance 0.75\,m}};
  \node[blk] (replan) [below=of compute]
    {REPLAN\\{\scriptsize NavFn recompute}};
  \node[blk] (evasion) [below=of drive]
    {LOCAL EVASION\\{\scriptsize local recovery attempt}};

  %% === FAIL-SAFE STATE ===
  \node[blk] (stop) [below=of replan]
    {FAIL-SAFE STOP\\{\scriptsize no safe trajectory}};

  %% === MAIN FLOW ===
  \draw[arr] (idle.east) -- (compute.west)
    node[midway, above=3pt, font=\scriptsize] {goal pose received};

  \draw[arr] (compute.east) -- (drive.west)
    node[midway, above, font=\scriptsize] {path valid};

  %% === TRACK -> LOCAL EVASION ===
  \draw[arr] (drive.south) -- (evasion.north);
  \node[font=\scriptsize, left] at ($(drive.south)!0.55!(evasion.north)+(-0.35,0)$)
    {path obstructed};

  %% === LOCAL EVASION -> REPLAN ===
  \draw[arr] (evasion.west) -- (replan.east);
  \node[font=\scriptsize, above] at ($(evasion.west)!0.5!(replan.east)+(0,0.12)$)
    {local evasion failed};

  %% === REPLAN -> COMPUTE PATH ===
  \draw[arr] (replan.north) -- (compute.south)
    node[midway, right, font=\scriptsize] {replan request};

  %% === HELPER COORDINATES ===
  \coordinate (lefttop) at ([xshift=-0.95cm]idle.west);
  \coordinate (leftbot) at ([xshift=-0.95cm]waypoint.west);
  \coordinate (topline) at ([yshift=1.00cm]drive.north);

  %% === TRACK -> WAYPOINT REACHED ===
  \draw[arr]
    (drive.north) -- (topline)
    -| (lefttop)
    -- (leftbot)
    -- (waypoint.west);

  \node[font=\scriptsize, above]
    at ($(topline)+(-4.9cm,0)$) {position error $< 0.75\,m$};

  %% === WAYPOINT REACHED -> IDLE ===
  \draw[arr]
    ([xshift=0.45cm]waypoint.north) -- ([xshift=0.45cm]idle.south);

  \node[font=\scriptsize, right]
    at ($([xshift=0.45cm]waypoint.north)!0.5!([xshift=0.45cm]idle.south)$)
    {next waypoint};

  %% === REPLAN -> FAIL-SAFE STOP ===
  \draw[dlink] (replan.south) -- (stop.north)
    node[midway, right, font=\scriptsize] {replan timeout 10\,s};

  %% === LOCAL EVASION -> TRACK TRAJECTORY ===
  \coordinate (pc1) at ([xshift=0.70cm]evasion.east);
  \coordinate (pc2) at ([xshift=0.70cm,yshift=1.75cm]evasion.east);

  \draw[dlink] (evasion.east) -- (pc1) -- (pc2) -- (drive.east);

  \node[font=\scriptsize, right]
    at ($(pc1)!0.52!(pc2)+(0.22,0)$) {local path clear};

\end{tikzpicture}}
  \end{minipage}

  \caption{Navigation overview: state machine and data flow}
  \label{fig:navigation_combined}
\end{figure}